% ---------------------------------------------------------------------
\chapter{Trecho do esquema JSON da ficha clínica}
\label{apendice:schema}

O esquema completo está no repositório do sistema, em \texttt{packages/ficha/schemas}. O \autoref{cod:schema} mostra parte de três seções da ficha inicial. A seção de história médica admite ditado. O questionário de dor admite apenas formulário, de modo que a IA não atribui notas de dor. O diagnóstico do Eixo~I é exclusivamente humano. As definições reutilizáveis, como a escala de 0 a 10, ficam no esquema comum e são referenciadas por \texttt{\$ref}.

\begin{lstlisting}[caption={Trecho do esquema da ficha inicial (versão 1)}, label=cod:schema]
{
  "$id": "urn:lord-uepg:ficha-inicial:v1",
  "type": "object",
  "additionalProperties": false,
  "properties": {
    "historia_medica": {
      "x-captacao": ["ditado", "formulario"],
      "type": "object",
      "additionalProperties": false,
      "properties": {
        "q3_alergias": {
          "title": "Histórico de alergias",
          "x-questao": "3",
          "x-obrigatorio": true,
          "$ref": "urn:lord-uepg:comum:v1#/$defs/texto"
        },
        "q4_usa_medicamento": {
          "title": "Faz uso de algum medicamento?",
          "x-questao": "4",
          "$ref": "urn:lord-uepg:comum:v1#/$defs/simNao"
        }
      }
    },
    "questionario_dtm": {
      "x-captacao": ["formulario"],
      "properties": {
        "q16_dor_atual": {
          "title": "Nota da dor na face neste momento (0 a 10)",
          "x-questao": "16",
          "$ref": "urn:lord-uepg:comum:v1#/$defs/escala0a10"
        }
      }
    },
    "diagnostico_eixo1": {
      "x-somente-humano": true,
      "properties": {
        "grupo1_muscular": {
          "enum": ["Ia_dor_miofascial",
                   "Ib_dor_miofascial_com_limitacao_abertura",
                   "sem_diagnostico", null]
        }
      }
    }
  }
}
\end{lstlisting}

O \autoref{cod:schema-ia} mostra o esquema derivado automaticamente para um campo da história médica, que é o formato enviado ao modelo de linguagem. A derivação resolve as referências, remove os metadados, mantém as instruções de preenchimento do próprio campo e exige, para cada campo, o valor e o trecho da transcrição que o justifica.

\begin{lstlisting}[caption={Esquema derivado para a IA (campo de alergias)}, label=cod:schema-ia]
"q3_alergias": {
  "type": "object",
  "additionalProperties": false,
  "description": "Histórico de alergias",
  "properties": {
    "valor": {
      "type": ["string", "null"],
      "maxLength": 4000,
      "description": "Registrar \"nega\" quando o paciente negar alergias."
    },
    "evidencia": {
      "type": ["string", "null"],
      "description": "Trecho literal da transcrição que justifica
                      o valor; null se não mencionado."
    }
  },
  "required": ["valor", "evidencia"]
}
\end{lstlisting}

% ---------------------------------------------------------------------
\chapter{Roteiro da avaliação de usabilidade}
\label{apendice:usabilidade}

\section{Tarefas}

\begin{enumerate}
	\item Acessar o sistema com o usuário de teste;
	\item Cadastrar o paciente do caso simulado;
	\item Registrar a anamnese e o exame clínico na condição sorteada (papel, formulário digital ou formulário com ditado e IA);
	\item Enviar a ficha para revisão;
	\item (Professores) Revisar, devolver com comentário e aprovar uma ficha.
\end{enumerate}

\section{Questionário SUS}

Cada item é respondido em escala de 1 (discordo totalmente) a 5 (concordo totalmente). O escore é obtido subtraindo 1 das respostas dos itens ímpares e subtraindo de 5 as respostas dos itens pares; a soma é multiplicada por 2,5, resultando em um valor de 0 a 100 \cite{brooke1996}.

\begin{enumerate}
	\item Eu acho que gostaria de usar este sistema com frequência.
	\item Eu achei o sistema desnecessariamente complexo.
	\item Eu achei o sistema fácil de usar.
	\item Eu acho que precisaria do apoio de um técnico para conseguir usar este sistema.
	\item Eu achei que as diversas funções deste sistema estavam bem integradas.
	\item Eu achei que havia muita inconsistência neste sistema.
	\item Eu imagino que a maioria das pessoas aprenderia a usar este sistema rapidamente.
	\item Eu achei o sistema muito complicado de usar.
	\item Eu me senti muito confiante usando o sistema.
	\item Eu precisei aprender muitas coisas antes de conseguir usar este sistema.
\end{enumerate}

\todo{substituir por uma tradução validada do SUS para o português e citar a fonte.}

\section{Entrevista semiestruturada}

\begin{enumerate}
	\item O que foi mais fácil e o que foi mais difícil em cada forma de registro?
	\item Você confiou nas sugestões da IA? Em que situações precisou corrigi-las?
	\item Em que momentos do atendimento você usaria o ditado? Em quais não usaria?
	\item O que falta no sistema para substituir a ficha em papel na rotina da Liga?
\end{enumerate}

% ---------------------------------------------------------------------
\chapter{Matriz de controle de acesso}
\label{apendice:rbac}

\begin{table}[htbp]
	\caption{Permissões por papel}
	\label{tb:rbac}
	\footnotesize
	\centering
	\begin{tabular}{|l|c|c|c|c|}
		\hline
		\textbf{Ação} & \textbf{Aluno} & \textbf{Professor} & \textbf{Pesquisador} & \textbf{Administrador} \\ \hline
		Cadastrar paciente & Sim & Sim & Não & Não \\ \hline
		Ver prontuário identificado & Próprios & Sim & Não & Não \\ \hline
		Preencher ficha (rascunho) & Sim & Sim & Não & Não \\ \hline
		Usar ditado e IA & Sim & Sim & Não & Não \\ \hline
		Aprovar ficha & Não & Sim & Não & Não \\ \hline
		Registrar adendo em ficha aprovada & Sim & Sim & Não & Não \\ \hline
		Ver painel do paciente & Próprios & Sim & Não & Não \\ \hline
		Ver painéis de coorte & Não & Sim & Sim & Não \\ \hline
		Exportar dados pseudonimizados & Não & Não & Sim & Não \\ \hline
		Gerenciar usuários & Não & Não & Não & Sim \\ \hline
		Consultar auditoria & Não & Não & Não & Sim \\ \hline
	\end{tabular}
	\fonte{Os autores (2026).}
\end{table}

\todo{confirmar a matriz com a coordenação da LORD.}
